\section{分析結果}
\label{sec:analysis}

本節推導三個解析結果，
刻畫 P2C 路由下慢故障緩解的根本限制與機制。

\subsection{操作包絡線}
\label{sec:envelope}

當偵測器將故障節點隔離時，剩餘健康節點的負載隨之上升。
直覺上，存在一個臨界點——隔離太多節點，
健康節點本身就會過載。以下定理量化了這個臨界點。

\begin{theorem}[操作包絡線]
\label{thm:envelope}
在具有負載因子 $L$ 的 $N$ 節點叢集中，
以隔離為手段的緩解策略，所能容忍的最大故障節點比例為：
\begin{equation}
  \fmax = 1 - \frac{L}{U}
\end{equation}
其中 $U$ 為使單節點排隊延遲達到不可接受程度的利用率上限。
\end{theorem}

\begin{proof}
當 $f$ 個節點被隔離（從活躍池中移除）時，
剩餘 $N - f$ 個健康節點必須承擔全部系統負載。
每節點的有效利用率變為：
\begin{equation}
  \rhoeff = \frac{\Lambda}{(N-f) \cdot \mu} = \frac{L \cdot N}{N - f}
\end{equation}
為使系統在可接受的排隊延遲下維持穩定，
需滿足 $\rhoeff \leq U$。對 $f$ 求解可得：
\begin{equation}
  f \leq N \left(1 - \frac{L}{U}\right)
\end{equation}
\end{proof}

簡言之，$L/U$ 是「最少需要多少比例的節點才能撐住負載」，
超出的部分才是可隔離的餘裕。
負載越高，餘裕越小。

\medskip

定理~\ref{thm:envelope}中的 $U$ 是一個自由參數——
它取決於對「可接受排隊延遲」的定義。
接下來從 M/D/1 模型推導其封閉形式，
如此一來，包絡線便成為無需調參的預測工具。
對於服務時間為確定值 $\tbase$ 的單伺服器佇列，
平均逗留時間為：
\begin{equation}
  T(\rho) = \tbase \cdot \frac{2 - \rho}{2(1 - \rho)}
  \label{eq:md1}
\end{equation}

\begin{corollary}[封閉式利用率上限]
\label{cor:U}
考慮一 M/D/1 佇列，其基底服務時間為 $\tbase$，SLO 目標以平均延遲衡量，則
利用率上限為：
\begin{equation}
  U = 1 - \frac{\tbase}{2 \cdot \text{SLO} - \tbase}
    = \frac{2(\text{SLO} - \tbase)}{2 \cdot \text{SLO} - \tbase}
  \label{eq:U}
\end{equation}
\end{corollary}

\begin{proof}
令公式~\ref{eq:md1}中 $T(U) = \text{SLO}$，
設 $\alpha = 2\,\text{SLO}/\tbase$，
直接求解一元一次方程式可得
$U = (\alpha - 2)/(\alpha - 1) = 1 - \tbase/(2 \cdot \text{SLO} - \tbase)$。
\end{proof}

以本文參數（$\text{SLO} = 50$\,ms，$\tbase = 10$\,ms）：
$U = 1 - 10/90 = 8/9 \approx 0.889$。全文使用 $U = 8/9$。

\medskip

\subsubsection*{推廣至 M/G/1 服務時間}
以下推廣並非主線論述所需，但可量化服務時間變異性對包絡線的影響。
當服務時間具有變異係數 $\Cs$ 時，
Pollaczek--Khinchine 公式給出
$T(\rho) = \tbase [1 + (1 + \Cs^2)\,\rho / (2(1-\rho))]$。
令 $T(U) = \text{SLO}$ 求解可得推廣上限：
\begin{equation}
  U(\Cs) = \frac{\alpha}{\alpha + 1 + \Cs^2}
  \quad\text{其中 } \alpha = 2\left(\frac{\text{SLO}}{\tbase} - 1\right)
  \label{eq:U_mg1}
\end{equation}

表~\ref{tab:Ucv}~展示服務時間變異性的影響。

\begin{table}[t]
\centering
\caption{不同服務時間分布下以平均延遲為基準的利用率上限 $U$
  與操作包絡線 $\fmax$（$L = 0.80$）。}
\label{tab:Ucv}
\begin{tabular}{@{}llccc@{}}
\toprule
分布 & $\Cs^2$ & $U$ & $\fmax$ & 最大故障數（$N{=}32$） \\
\midrule
M/D/1 & 0 & 0.889 & 10.0\% & 3.2 \\
M/G/1（溫和） & 0.25 & 0.865 & 7.5\% & 2.4 \\
M/M/1 & 1.0 & 0.800 & 0.0\% & 0.0 \\
\bottomrule
\end{tabular}
\end{table}

M/M/1（$\Cs^2 = 1$）的結果值得特別關注：
$U = 0.800$，意味著在 $L = 0.80$ 時故障容忍度為\emph{零}——
任何慢故障都會立即違反 SLO。
此結果量化了一個直覺認識——服務時間變異性會侵蝕系統餘裕，
也表明確定性 M/D/1 假設提供了故障容忍度的\emph{樂觀}上界。

\begin{remark}[平均值上限的有效性與緊密性]
\label{rem:mean-p99}
表~\ref{tab:Ucv}~由平均延遲導出 $U$，
但 SLO 通常定義於 P99。
對 M/D/1（$\Cs^2 = 0$），差距很小：
確定性服務時間使回應時間的變異很小，
故 $U_{\text{P99}} \gtrsim U_{\text{mean}}$，
公式為\emph{保守的}（安全的）。
對 M/M/1（$\Cs^2 = 1$），差距很大：
模擬顯示 P99 在 $L \approx 0.70$ 即超越 50\,ms SLO，
基於平均值的 $U$ 對高變異分布而言是\emph{樂觀的}。
就緊密性而言，
在 $L = 0.80$、$N = 32$ 時理論上限為 3.2 個節點，
而實驗中隔離 4 個節點仍可達成 P99 $< 50$\,ms，
此餘裕源自平均延遲與 P99 之間的差距。
包絡線略為保守，但預測的通過／失敗邊界與實驗相差不超過一至兩個節點。
\end{remark}

\medskip

\subsubsection*{非隔離策略}
定理~\ref{thm:envelope}~中的上界假設故障節點被完全隔離。
若策略僅\emph{削減}故障節點流量而非完全隔離，
則故障節點貢獻 $\mu/s_i$ 的有效服務率，
有效利用率為：
\begin{equation}
  \rhoeff = \frac{L \cdot N}{(N-f) + \sum_{i \in F} 1/s_i}
\end{equation}
其中 $F$ 為故障節點集合。
此值始終低於完全隔離的情況，
證實操作包絡線為\emph{保守估計}。

以 $N = 32$、永久 $5\times$ 故障的配置，
在四個負載等級下驗證定理~\ref{thm:envelope}
（表~\ref{tab:envelope}）。

\begin{table}[t]
\centering
\caption{操作包絡線驗證。理論預測最大故障比例；
  實驗欄位顯示自適應策略是否滿足 P99 $< 50$\,ms 之 SLO。}
\label{tab:envelope}
\begin{tabular}{@{}ccccl@{}}
\toprule
$L$ & $\fmax$ & $f_{\max}$ 節點數 & 實驗邊界 & 吻合 \\
\midrule
0.70 & 21.3\% & 6.8 & 8 通過, 10 失敗 & \checkmark$^\dagger$ \\
0.80 & 10.0\% & 3.2 & 4 通過, 6 失敗 & \checkmark \\
0.85 & 4.4\% & 1.4 & 2 通過, 4 失敗 & \checkmark$^\dagger$ \\
0.90 & $<0$ & --- & 0 通過, 2 失敗 & \checkmark \\
\bottomrule
\end{tabular}
\end{table}

該公式無需任何擬合參數即可預測通過／失敗邊界。
標記 $\dagger$ 的列中，實驗通過的故障數超過理論上限
（如 $L=0.70$ 時理論上限為 6.8 個節點，但 8 個仍通過），
這是因為 $U$ 由\emph{平均}延遲導出，而 P99 的有效上限略高
（見備註~\ref{rem:mean-p99}），使包絡線成為保守估計。
$L = 0.90$ 超過 $U = 8/9$，即使無故障亦無法滿足 SLO。

圖~\ref{fig:envelope}~以圖形呈現各負載等級下的操作包絡線，
展示安全區與過載區之間的明確邊界。

\begin{figure}[t]
\centering
\includegraphics[width=0.85\linewidth]{operating_envelope.png}
\caption{操作包絡線：理論與實驗對照。
  邊界 $\fmax = 1 - L/U$（虛線）
  劃分安全區（P99 $< 50$\,ms）與過載區，
  驗證過程無需任何擬合參數。}
\label{fig:envelope}
\end{figure}

\subsection{嚴重度非單調性}
\label{sec:nonmono}

\begin{assumption}[平均場 P2C 近似]
\label{asm:meanfield}
在 P2C 加權隨機抽樣下，穩態時各節點所接收的流量
與其有效服務容量成正比。
即容量為 $\mu_i$ 的節點接收總流量的
$\mu_i / \sum_j \mu_j$ 份額。
\end{assumption}

當節點數量夠大且抽樣按容量加權時，此近似趨於精確。
對有限叢集而言，P2C 的佇列深度比較
會使故障節點的流量分配產生輕微的正偏差
（見表~\ref{tab:pfault}：$f=4$ 時低估 10--14\%，
$f=8$ 時 $< 1$\%）。

\medskip

\begin{theorem}[臨界嚴重度]
\label{thm:sstar}
在假設~\ref{asm:meanfield}~下，
流量到達 $f$ 個（減速因子為 $s$）故障節點的比例為：
\begin{equation}
  \pfault(f, s, N) = \frac{f}{s \cdot (N-f) + f}
  \label{eq:pfault}
\end{equation}
使故障節點的延遲影響在第 $q$ 百分位上不可見的臨界嚴重度為：
\begin{equation}
  \sstar = \frac{f \cdot q}{(1-q)(N-f)}
  \label{eq:sstar_deriv}
\end{equation}
\end{theorem}

換言之：P2C 路由會讓越嚴重的故障節點分到越少的流量。
當嚴重度夠高時，故障流量的佔比降到 $1-q$ 以下——
此時故障請求在第 $q$ 百分位上統計不可見。
$\sstar$ 就是這個臨界點。

\begin{proof}
在 P2C 下，流量按有效服務容量的比例分配。
健康節點容量為 $\mu$；故障節點容量為 $\mu/s$。
故障節點的流量佔比化簡為
$\pfault = f / [s(N{-}f) + f]$（公式~\ref{eq:pfault}）。
系統延遲服從混合分布，
當 $\pfault > 1-q$ 時故障請求在第 $q$ 百分位可見。
令 $\pfault = 1-q$ 並對 $s$ 求解即得 $\sstar$。
\end{proof}

\medskip

\subsubsection*{非單調性機制}
當 $s < \sstar$ 時，$\pfault > 1-q$，
故障節點的延遲主導第 $q$ 百分位。
系統 $P_q$ 隨 $s$ 增大而升高，
因故障節點的服務時間為 $s \cdot \tbase$。
當 $s > \sstar$ 時，$\pfault < 1-q$，
故障請求在第 $q$ 百分位統計上不可見。
系統 $P_q$ 回落至健康節點值。
因此，\emph{最差}的 $P_q$ 發生在 $s = \sstar$ 附近，
而非 $s \to \infty$。
此即嚴重度非單調性：
將故障嚴重度提高至 $\sstar$ 以上，
尾延遲反而\emph{降低}。

換言之，系統管理員面臨的最棘手情境並非極端故障（$s \gg \sstar$），
而是中度故障（$s \approx \sstar$）——這恰恰是最難偵測的區間。

\subsubsection*{百分位相關風險}
同一故障在不同百分位上的可見性可能截然不同，此即百分位相關風險。
對 $f = 4$，$N = 32$：$s^{*}_{P99} = 14.1$ 但
$s^{*}_{P999} = 142.7$。
$20\times$ 減速對 P99 不可見（$20 > 14.1$）
但對 P999 可見（$20 < 142.7$），
這解釋了實驗中觀察到的現象：
P99 $= 49$\,ms（正常）但 P999 $= 17{,}427$\,ms（災難性）。

表~\ref{tab:pfault}~驗證上述流量分配模型。

\begin{table}[t]
\centering
\caption{到達故障節點的流量比例：理論與實驗對照
  （無緩解，$N=32$，$L=0.80$）。}
\label{tab:pfault}
\begin{tabular}{@{}ccccc@{}}
\toprule
$f$ & $s$ & 理論 $\pfault$ & 實驗值 & 比率 \\
\midrule
4 & $3\times$ & 4.55\% & 5.20\% & 1.14 \\
4 & $5\times$ & 2.78\% & 3.12\% & 1.12 \\
4 & $10\times$ & 1.41\% & 1.56\% & 1.11 \\
4 & $20\times$ & 0.71\% & 0.78\% & 1.10 \\
\midrule
8 & $3\times$ & 10.00\% & 10.42\% & 1.04 \\
8 & $5\times$ & 6.25\% & 6.26\% & 1.00 \\
8 & $10\times$ & 3.23\% & 3.23\% & 1.00 \\
8 & $20\times$ & 1.64\% & 1.64\% & 1.00 \\
\bottomrule
\end{tabular}
\end{table}

模型在 $f=4$ 時系統性低估 10--14\%
（因 P2C 的佇列深度比較機制使較慢節點接收到略多於模型預測的流量），
而在 $f=8$ 時則精確至三位有效數字。
$f = 4$ 時的系統性正偏差（比率 $\approx 1.10$--$1.14$）
源自 P2C 佇列深度比較中的平手打破機制，此殘餘偏差以 $O(1/N)$ 的速率遞減。
就實際用途而言，平均場近似足以預測臨界嚴重度 $\sstar$
及 P99 可見性邊界。

圖~\ref{fig:heatmap}~展示不同故障數量及減速因子組合下的
嚴重度非單調性。
熱力圖清楚顯示中度嚴重度
（如 4 個故障 $10\times$：P99 $= 439$\,ms）
產生的尾延遲比極端嚴重度
（如 4 個故障 $20\times$：P99 $= 49$\,ms）更差。

\begin{figure}[t]
\centering
\includegraphics[width=0.85\linewidth]{severity_heatmap.png}
\caption{嚴重度非單調性熱力圖。$L = 0.80$，$N = 32$ 下，
  按故障數量與減速因子呈現 P99 延遲（ms）。
  非單調模式清晰可見：在每一種故障數量下，
  存在一個臨界嚴重度，超過該值後 P99 急遽下降。}
\label{fig:heatmap}
\end{figure}

\subsection{SHED--P2C 互補性}
\label{sec:complementarity}

操作包絡線與臨界嚴重度界定了緩解的可能範圍。
接下來分析 P2C 與顯式權重削減（SHED）之間的互動，
證明兩者為互補機制——而非冗餘。

\subsubsection*{P2C 下的兩階段路由}
P2C 路由做出兩個依序的決策：
\begin{enumerate}
\item \emph{候選者選取}：依權重 $w_i$ 的比例抽取兩個工作節點。
  故障節點被選為候選者的機率約為 $2w_i / \sum_j w_j$。
\item \emph{佇列深度比較}：在兩個候選者中，
  將請求路由至佇列較短的一方。
\end{enumerate}
SHED 作用於第一階段，降低故障節點的 $w_i$；
P2C 的隱式迴避則作用於第二階段，透過佇列深度比較實現。
因此，第二階段只能在第一階段所選定的候選者之間做決策。

\begin{proposition}[SHED--P2C 互補性]
\label{thm:complementarity}
在採用 $d = 2$ 選擇的 P2C 路由下，
考慮一個 $N$ 節點叢集中的單一故障節點（嚴重度為 $s$），
令 $\Delta_{\text{SHED}}(N, s)$ 為對其施加 SHED 權重削減後的 P99 改善量。則：
\begin{enumerate}
\item[(a)] \emph{規模衰減}：$\Delta_{\text{SHED}}(N, s) = O(N^{-\alpha})$，
  其中 $\alpha \approx 1.3$（中度嚴重度區間）。
\item[(b)] \emph{嚴重度區間}：$\Delta_{\text{SHED}}$ 在中等 $s$
  （通常 $3$--$5\times$）時最大化，
  在 $s < 1.5$ 及 $s > \sstar$ 時可忽略。
\item[(c)] \emph{實務等價}：在 P2C 下，最低權重 $w_{\min} = 0.05$
  可達到與完全隔離（$w = 0$）相差 $O(w_{\min}/N)$ 以內的 P99。
\end{enumerate}
\end{proposition}

直覺上：P2C 的佇列比較已是一道強大的防線。
SHED 在此之上再加一層過濾，
但因為第一道防線已經很有效，
第二道的邊際效益隨叢集規模衰減。

\medskip

此結論源自兩階段結構。無 SHED（均勻權重）時，
每次請求選到故障節點作為候選者的機率約為 $2/N$。
在 SHED 權重 $w_{\min}$ 下，此機率降為 $\approx 2w_{\min}/N$——
候選率降低了 $1/w_{\min}$ 倍。
然而，P2C 的第二階段透過佇列比較已在多數時候拒絕故障節點。
到達故障節點的\emph{淨}流量在無 SHED 時為 $O(1/N)$，
有 SHED 時為 $O(w_{\min}/N)$。
隨著 $N$ 增長，故障節點的基準流量已可忽略，
SHED 的邊際效益因而消失。

\subsubsection*{容量損失邊界}
上述互補性有一個重要的例外：
當健康節點的有效負載
$L_{\text{healthy}} = L \cdot N / (N - f_{\text{shed}})$
超過排隊膝點（$\approx 0.88$--$0.90$）時，
SHED 的權重削減反而有害——
重新分配的流量使健康節點飽和，
導致的排隊延遲增加超過原先的改善。

表~\ref{tab:shed_complementarity}~於不同叢集大小、
嚴重度等級與故障數下驗證上述互補性命題。

\begin{table}[t]
\centering
\caption{SHED 在 P2C 下相對於無緩解的 P99 改善（ms）。
  SHED 將權重降至 $w = \max(0.05, 1-\sigma)$；ISO 設 $w = 0$。
  所有數值為 10 次執行的平均值。}
\label{tab:shed_complementarity}
\begin{tabular}{@{}lccccl@{}}
\toprule
配置 & no\_mit & SHED & ISO & SHED $\Delta$ & 備註 \\
\midrule
\multicolumn{6}{@{}l}{\emph{規模掃描（$s = 3\times$，$L = 0.70$）}} \\
$N=8$  & 79.7 & 34.6 & 34.3 & $-45.1$ & 小 $N$ 時效益大 \\
$N=16$ & 59.3 & 29.8 & 28.2 & $-29.5$ & \\
$N=32$ & 46.6 & 27.8 & 26.9 & $-18.8$ & \\
$N=64$ & 28.5 & 26.8 & 26.4 & $-1.7$  & 大 $N$ 時可忽略 \\
\midrule
\multicolumn{6}{@{}l}{\emph{嚴重度掃描（$N = 16$，$L = 0.70$）}} \\
$s=1.5$ & 29.3 & 28.0 & 28.1 & $-1.3$ & 低於偵測閾值 \\
$s=3$   & 59.3 & 29.8 & 28.2 & $-29.5$ & 效益峰值 \\
$s=5$   & 99.7 & 29.7 & 28.2 & $-70.0$ & 接近 $\sstar$ \\
$s=10$  & 34.6 & 34.6 & 34.6 & $+0.0$ & 高於 $\sstar$：已自行隔離 \\
\midrule
\multicolumn{6}{@{}l}{\emph{無佇列感知路由（加權隨機）}} \\
$s=3$, WR & 17757 & 2060 & --- & $-15697$ & 無 P2C 時 SHED 至關重要 \\
\bottomrule
\end{tabular}
\end{table}

最後一列表明 SHED 本身\emph{並非}無效——
在無佇列感知路由下可帶來 15,697\,ms 的改善。
P2C 的佇列深度比較大幅壓縮了 SHED 的邊際貢獻。
此結果證實兩種機制互補：P2C 提供主要防禦，
SHED 在候選者選取階段提供次級過濾。

\medskip

在 P2C 下，SHED 與 ISO 的 P99 差距始終 $< 2$\,ms
（表~\ref{tab:shed_complementarity}），
而無緩解與任何主動策略的差距則為 5--70\,ms。
此一數量級的差異意味著，
工程上的首要挑戰是\emph{多快}偵測到故障節點，
而非\emph{多精確}地調降其權重。
在 P2C 下，粗略的二元偵測器（「降權或不降權」）
即可達到近乎最優的效果。

\subsection{對緩解設計的啟示}

上述公式將任意 $(f, s, L, N)$ 配置劃分為三個區域：

\begin{enumerate}
\item 不可行（$f/N > \fmax$）：
  任何策略都無法維持 SLO。容量理論極限不可逾越。
\item P2C 充分（$f/N \leq \fmax$ 且 $s > \sstar$）：
  P2C 透過隱式佇列深度迴避自動處理故障節點。
  顯式緩解非必要，反而可能引入額外開銷。
\item 需要緩解（$f/N \leq \fmax$ 且 $s \leq \sstar$）：
  故障比例在容量限制內，但 P2C 無法充分避開故障節點。
  需要主動偵測與流量重導。
  在 P2C 下，漸進式權重削減與完全隔離的差異可忽略
  （$< 2$\,ms）——關鍵在於偵測速度。
\end{enumerate}

結合互補性命題（第~\ref{sec:complementarity}~節），
此框架釐清顯式緩解的適用條件：
在區域 2，P2C 自行隔離故障節點，
無需偵測或緩解——但監控 P999 至關重要，
因高嚴重度故障在更深的百分位仍然可見。
在區域 1，任何策略皆無濟於事。
